\textit{Can video pretraining serve as a foundation for embodied navigation?} We investigate this question in a synthetic environment that simplifies the world into a 2D grid and the rules of physics to those in Conway's Game of Life, where transitions are known exactly. At the core of our experiments is the \textit{Frame Crossing Game} in which an agent must navigate across a $96 \times 96$ grid without hitting the live cells, anticipating the evolution of surrounding cells governed by the game of life rules.  We find that video pretraining provides a consistent advantages over reinforcement learning from random initialization, but the nature and magnitude of this benefit varies with environment complexity. In the simplest environments, pretraining accelerates convergence: both pre-trained and trained-from-scratch agents ultimately reach the same level of progress, but the pretrained agent does so substantially faster. In the most difficult environments, no meaningful difference emerges between the two, as all models struggle to make measurable progress regardless of initialization. The most striking advantage appears at intermediate difficulty levels: here, the pretrained agent not only converges faster but converges to a higher level of progress.
